\section{Resumo e Atividade Prática}

\begin{frame}{Tabela Resumo dos Tipos de Funções}
    Aqui está o mapa mental para você nunca mais esquecer os tipos de funções no Python:
    
    \vspace{0.4cm}
    \begin{table}
        \centering
        \begin{tabular}{lccl}
            \toprule
            \textbf{Tipo de Função} & \textbf{Entrada?} & \textbf{Saída?} & \textbf{Exemplo Prático (Redes)} \\
            \midrule
            Sem Param. e Sem Ret. & Não & Não & \texttt{exibir\_banner()} \\
            Com Param. e Sem Ret. & Sim & Não & \texttt{ping("10.0.0.1")} \\
            Sem Param. e Com Ret. & Não & Sim & \texttt{obter\_ip\_loopback()} \\
            Com Param. e Com Ret. & Sim & Sim & \texttt{verificar\_ip\_privado(ip)} \\
            \bottomrule
        \end{tabular}
        \caption{Comparação dos quatro tipos de funções em Python.}
    \end{table}
\end{frame}

\begin{frame}{Por que usar Funções? (Boas Práticas)}
    Como técnicos em redes, sabemos que a organização evita desastres (como cabos espalhados no rack!). No código, é a mesma coisa:
    
    \vspace{0.3cm}
    \begin{itemize}
        \item \textbf{Evita Duplicação (Princípio DRY - Don't Repeat Yourself):} Escreva uma vez, use quantas vezes quiser.
        \item \textbf{Facilidade de Manutenção:} Se o cálculo da sub-rede mudar, você altera apenas dentro da função, e não em centenas de linhas de código.
        \item \textbf{Organização e Legibilidade:} O código fica estruturado em módulos fáceis de ler, parecendo comandos do sistema operacional.
    \end{itemize}
\end{frame}

\begin{frame}[fragile]{Laboratório de Redes: Mini Desafio}
    \textbf{Exercício de Fixação:} Crie um script Python que automatize a verificação de uma lista de IPs da rede escolar.
    
    \begin{lstlisting}[language=Python]
# 1. Crie uma funcao que recebe um IP e retorna se ele esta ativo
# (Para o teste, retorne True se o IP terminar com par e False se impar)
def testar_host(ip):
    ultimo_octeto = int(ip.split(".")[-1])
    if ultimo_octeto % 2 == 0:
        return True
    return False

# 2. Varra uma lista de IPs e use a funcao criada para exibir o status
ips_rede = ["192.168.1.10", "192.168.1.15", "192.168.1.20"]

for ip in ips_rede:
    if testar_host(ip):
        print(f"Host {ip} -> ONLINE")
    else:
        print(f"Host {ip} -> OFFLINE")
    \end{lstlisting}
\end{frame}

\begin{frame}{Conclusão: Próximos Passos}
    \begin{center}
        \LARGE \textbf{Parabéns! Você deu o primeiro passo na Automação de Redes!} \\
        \vspace{0.5cm}
        \normalsize Na próxima aula de Programação de Computadores, vamos integrar essas funções com a biblioteca \texttt{socket} para interagir com roteadores de verdade na nossa rede local! \\
        \vspace{0.8cm}
        \textbf{Dúvidas?} \\
        \texttt{reginaldo.fernandes@ifce.edu.br}
    \end{center}
\end{frame}
